
\title{Chemical Reactions in H-N-Silicate Systems at High P--T: Implications for nitrogen partitioning between envelope and interior of sub-Neptunes}
\author{Spoke person: Taehyun Kim}
\maketitle

%%%%%%%%%%%%%%%%%%%%%%%%%%%%%%%%%%%%%%%%%%%%%%%%%%%%%%%%%%
\section{Abstract ($<2000$ character)}

%\begin{cbox}
%The abstract should be written as goal $\rightarrow$ introduction $\rightarrow$ detailed plan.  I changed the order of sentences to fit this flow.  Please use this flow next time.
%\end{cbox}


% 2114
We request beamtime to investigate reaction between olivine/feropericlase and N-H fluids at high P-T to understand partitioning of nitrogen between the envelope and the interior of sub-Neptune exoplanets. 

Recent JWST measurements have identified NH$_3$, HCN, and N$_2$ in hydrogen-rich atmosphere of exoplanets \cite{heng_ANALYTICAL_2016, bean_Nature_2021, radecka_Nitrogen_2025}. 
While possible ingassing of nitrogen under this condition can impact the interpretation of the observation and the modeling of sub-Neptunes, partitioning of nitrogen between the atmosphere and the interior under highly reducing conditions expected in hydrogen-rich environment of sub-Neptunes remains poorly understood. 

Our recent experiments at 13IDD have revealed that H$_2$/H$_2$O fluids can react extensively with silicates and oxides, producing metallic alloys, hydrides, SiH$_4$, and hydrous phases, revealing large exchange of elements between the envelope and the interior \cite{nisr_Large_2020, kim_Atomicscale_2021, kim_Stability_2023, horn_Building_2025}. 
These studies have also demonstrated that synchrotron high-pressure experiments can provide critical data for understanding exoplanets and early terrestrial planets (like Earth).

Our preliminary experiments at 13IDD beamline showed that N-H fluids can react with silicates and oxides to form nitrides and Fe-N alloys at 10~GPa and over 3000~K (Fig.~\ref{Olivine-N2-H2}). 
At lower temperatures, retainment of N-H bonding in a fluid phase was found. 
These observations suggest that nitrogen can be distributed between the envelope and the interior through fluid-silicate reactions. 


In the proposed beamtimes at 13IDD, we will investigate the reaction and its pressure and temperature-dependent variations using in situ X-ray diffraction in a laser-heated diamond-anvil cell at {3--30~GPa and 1500--4000~K} conditions relevant to the envelope-interior boundary of sub-Neptune exoplanets. 
After the synchrotron experiments, Raman spectroscopy and electron microscopy will be conducted, which will provide additional data on reaction products and phase distributions. 
The resulting dataset will provide {silicate stability under N-H fluid,} nitride formation, Fe-N alloy formation, and the N-H bearing molecule formation at different P-T conditions, providing key (yet currently missing) data for understanding size and age dependent changes in the chemical interaction.
Therefore, the proposed data will advance our understanding on atmosphere-interior interactions in early terrestrial planets and sub-Neptune exoplanets.

%%%%%%%%%%%%%%%%%%%%%%%%%%%%%%%%%%%%%%%%%%%%%%%%%%%%%%%%%%%%%%%%%%%%
\section{What is the scientific or technical purpose and importance of the proposed research?  ($<2000$ character)}

% 1835
Sub-Neptunes are among the most abundant classes of exoplanets discovered to date. 
Their low bulk densities (2--4~g/cm$^3$) can be explained by envelopes rich in either hydrogen or water overlying silicate and metallic interiors (gas dwarfs and water worlds, respectively) \cite{bean_Nature_2021}.
The conditions and compositions in these settings are dramatically different from those of terrestrial planets, requiring new data to understand the astrophysical data of the exoplanet class.

Recent high P-T studies have demonstrated that envelope materials can react with materials of the interior under the conditions, producing metal alloys, hydrides, and H$_2$O. 
Such reaction can modify the compositions of both atmosphere and interior \cite{horn_Building_2025, kim_Stability_2023, miozzi_Experiments_2025}. 
However, previous studies focused primarily on pure hydrogen fluid interacting with silicates, whereas astrophysical surveys have revealed much more complex atmosphere compositions of exoplanets, including abundant presence of nitrogen \cite{radecka_Nitrogen_2025}.

To our knowledge, interaction between N-H fluids and silicates have not been systematically investigated at relevant conditions for the planets. %\sidecomment{Note that by writing N$_2$-H$_2$ you are assuming that nitrogen and hydrogen will exist in these specific forms, not the other forms, such as NH3 or even ingassed forms.}
It is unknown whether nitrogen remains in the fluid phase of the envelope, or forms nitrides, Fe-N alloys, or dissolves into silicates (therefore ingassed) under the highly reducing conditions expected in hydrogen-rich environment of sub-Neptunes.

Our preliminary experiments suggest that nitrogen can be incorporated into planetary interiors (Fig.~\ref{Olivine-N2-H2}). 
The formation of nitrides and Fe-N alloys implies strong nitrogen exchange between the envelope and the interior, affecting evolution and interpretation of the measured atmosphere spectra of these planets.

The primary objective of our proposed experiments is to understand nitrogen partitioning between fluid and silicates at high P-T.
In our recent studies \cite{horn_Building_2025}, we also found that the experiment designed and executed for exoplanets can also provide new insights on early evolution of terrestrial planets (including Earth) \cite{young_Earth_2023}. %\sidecomment{cite Young 2023 Nature paper here.}
Considering importance of nitrogen for Earth's atmosphere, our new data can potentially impact our understanding of how the atmosphere chemistry can be affected by underlying geology in early terrestrial planets.


%%%%%%%%%%%%%%%%%%%%%%%%%%%%%%%%%%%%%%%%%%%%%%%%%%%%%%%%%%%%%%
\section{Why do you need the APS for this research?  ($<2000$ character)}

% 983
The proposed chemical system will produce low-Z materials with weak X-ray scattering factors. 
The reaction products include nitrides, hydrides, and silicates.
Furthermore, the amount of the materials will be extremely small under high-pressure and high-temperature conditions in a diamond-anvil cell. 
To overcome these factors, the high brilliance X-ray beam available at APS is critical for reliable phase identification.
The small apertures of diamond-anvil cells require high energy X-ray beam for ensuring detection of low d-spacing critical for accurate identification of phases and determination of unit-cell volumes. 


Our experimental plan also includes very short heating durations (1-2 seconds). 
The reason is that during longer laser heating in fluid, silicate melts can migrate away from the heated area.
Therefore, short heating duration is required.
To monitor reactions during these short heating events, X-ray beam must be sufficiently intense to produce high-quality diffraction patterns of weakly scattering nitride, hydride, and silicate phases. 
The laser-heated region is typically less than 20 $\mu$m in diameter at high pressure, necessitating a highly collimated X-ray beam with a size of less than 2 $\mu$m.

%%%%%%%%%%%%%%%%%%%%%%%%%%%%%%%%%%%%%%%%%%%%%%%%
\section{Describe why you are choosing your requested beamline ($<2000$ character)}

% 1435
Beamline 13IDD is designed and optimized for high-pressure experiments. 
The beamline provides high-brilliance X-rays with focusing optics, coupled with state-of-the-art experimental instrumentation for in situ laser heating in a diamond-anvil cell. 
The laser heating system at 13IDD is a crucial component for conducting our proposed experiments where silicates will be heated in hydrogen-rich fluids to temperatures approaching 4000 K.

The beamline features pulsed-laser heating mode and burst (ramp) heating mode, essential for achieving sufficiently high temperatures for melting all the materials including silicates and fluid medium at high pressures. 
These two heating methods are particularly useful for reaching temperatures above 3000 K, where continuous laser heating could induce severe fluid convection problems, thereby lowering laser coupling efficiency. 
In other words, these dedicated heating techniques enable us to obtain data on chemical reactions without suffering from dynamic instability in the sample chamber during melting.

The beamline also provides an extremely well-collimated X-ray beam on the scale of 1-2 micrometers, and the X-ray beam size can be further reduced to a few hundred nanometers. 
This smaller beam size helps to minimize the impact from the radial thermal gradient in the heated area. 
Additionally, the user-friendly interface of the software programs allow users to obtain and analyze high-quality XRD images and temperature data efficiently during beamtime.



%%%%%%%%%%%%%%%%%%%%%%%%%%%%%%%%%%%%%%%%%%%%%%%%
\section{How many visits during the proposal lifespan do you expect to need? ($<2000$ character)} 

%2093
We request 12 shifts for this project, with 3 shifts allocated per visit. 
We plan to prepare four to five diamond-anvil cells (DACs) per visit, depending on the number of shifts allocated.

For the first visit, we will prepare two DACs containing ferropericlase in a N$_2$-67\%H$_2$ medium and two DACs containing olivine in a N$_2$-67\%H$_2$ medium. 
A gas loading system at ASU is capable of loading gas mixtures and we will prepare these samples before the beamtime trip \cite{fu_Core_2023}.
Based on our previous experience for similar experiments at 13IDD, approximately 5 hours are required per DAC, including sample alignment, laser heating, and post-heating 2D XRD mapping (5 hours per DAC and 4 DACs, totaling 20 hours). 
The target pressures for the first visit will be 5--14 GPa, with temperatures ranging from 1500 to 4000 K. 
An additional 4 hours will be required to collect 2D XRD maps of the recovered samples at ambient pressure.

During the second visit, we will investigate higher-pressure conditions (15-30 GPa). 
We will prepare two DACs containing ferropericlase in a N$_2$-67\%H$_2$ medium and two DACs containing olivine in a N$_2$-67\%H$_2$ medium. 

\added{During the third visit, we will target very high temperature (4000~K) at very low pressures (2-5 GPa), which is important for understanding smaller young sub-Neptune exoplanets.
The same two compositions we discussed above will be investigated.}
Because pressure instabilities may occur during laser heating to extremely high temperatures at low pressures, we will prepare one or two more backup DACs. 
If needed, we will further reduce heating durations. 

During the fourth visit, we will focus on verifying the reproducibility of the results obtained in first three visits and completing the dataset sufficient for understanding pressure- and temperature-dependent variation in chemical reactions in the N-H-silicate system. 
If previously unknown phases are observed in first three visits, we will synthesize them and measure their equations of state in this beamtime.


%%%%%%%%%%%%%%%%%%%%%%%%%%%%%%%%%%%%%%%%%%%%%%%%
\section{Describe/provide a list of samples ($<2000$ character)} 

% 390

We will use natural San Carlos olivine ((Mg$_{0.9}$Fe$_{0.1}$)$_2$SiO$_4$), which have been also used in our previous experiments \cite{kim_Atomicscale_2021, kim_Solubility_2022}. 
For ferropericlase, we will use (Mg$_{0.75}$Fe$_{0.25}$)O synthesized in a multi-anvil press at FORCE, ASU. 
Combining these two compositions, we will be able to isolate the effects of Si in the chemical reaction and cover a wide range of Mg/Si ratios for the limited beamtimes.

\replaced{Nitrogen and hydrogen will be mixed in 1:3 ratio in our gas loading system at ASU.
The gas mixture will be captured together with the silicate/oxide samples to 1500~bars in DACs before the beamtime trips.}{A N$_2$-H$_2$ mixed gas (33 mol\% N$_2$ and 67 mol\% H$_2$) will be used as the reactive fluid during high-pressure and high-temperature experiments.}



%%%%%%%%%%%%%%%%%%%%%%%%%%%%%%%%%%%%%%%%%%%%%%%%
\section{Overview of the experimental plan and procedures, including sample usage  ($<2000$ character)} 


We will use diamond anvil culet diameters of 300 and 400 $\mu$m for laser-heating experiments at 15-30 and 5-14 GPa, respectively. 
A thin foil ($\sim$10 $\mu$m thick) of compressed olivine or ferropericlase powder will be loaded into a drilled gasket hole (typically 200-300 $\mu$m in diameter) together with a N$_2$-H$_2$ gas mixture as a medium. 
Sample spacer grains will be used to support the foil and prevent direct contact with the diamond anvils. 
The gases will be loaded using a gas-loading system at ASU.

During heating to temperatures over $\sim$4000 K at lower pressures (e.g., 5 GPa), mechanical relaxation of the DACs can occur, which can result in loss of pressure and therefore leakage of the pressure medium. 
Such failures can occur primarily during long heating. 
Therefore, we will reduce heating durations and use smaller gasket holes to reduce the risk of potential sample loss during extreme high temperature heating.

Ruby fluorescence will be used for measuring pressure. 
We will employ both burst-mode and continuous laser-heating methods and evaluate their relative effectiveness based on the experimental results during our first visit. 
An EIGER2 X CdTe 9M detector at the beamline will enable rapid acquisition of high-quality diffraction patterns during and after heating, with particular emphasis on identifying nitrides, Fe-N alloys, hydrides, and other reaction products formed during heating. 
Two-dimensional XRD maps will be collected after heating at high pressures and after decompression to ambient pressure.

Following the laser heating, Raman spectra will be measured using the GSECARS Raman system to characterize reaction products both at high pressure and after recovery to ambient conditions. 
To maximize efficiency during beamtime, one group member will conduct Raman measurements while the other researcher performs laser-heating experiments on another DAC. 
After beamtime, the recovered samples will be characterized using an FIB-SEM-EDS instrument at ASU to determine composition and microstructure, providing additional information.

%%%%%%%%%%%%%%%%%%%%%%%%%%%%%%%%%%%%%%%%%%%%%%%%
\section{Describe the team's previous experimental experience with synchrotron radiation ($<2000$ character)}

Taehyun Kim, the principal investigator (PI), has conducted LHDAC experiments at 13IDD at GSECARS, 16IDB at HPCAT, and P02.2 at PETRA III, and has performed numerous XRD experiments at ALS, SSRL, and PLS-II.

Sang-Heon Dan Shim, the PI's advisor, has conducted numerous XRD experiments at APS, CHESS, NSLS, SSRL, and ALS over the past 30 years.

Allison Pease and Alem Tukic, members of Shim's group, have carried out synchrotron X-ray diffraction experiments at APS, including at 13IDD.

Taehyun Kim and Sang-Heon Dan Shim have already conducted similar experiments as proposed here in recent years at the 13IDD beamline \cite{kim_Stability_2023, horn_Building_2025}.



%%%%%%%%%%%%%%%%%%%%%%%%%%%%%%%%%%%%%%%%%%%%%%%%
\section{List publications resulting from work done at the APS  ($<2000$ character)}

%\inlinecomment{Please identify the beamline(s) where the work was done.}

Selected recent publication based on past APS work includes:

Fu, S., Chariton, S., Prakapenka, V. B., and Shim, S.-H. (2023). Core origin of seismic velocity anomalies at Earth's core-mantle boundary. Nature. https://doi.org/10.1038/s41586-023-05713-5. 

Kim, T., O'Rourke, J. G., Lee, J., Chariton, S., Prakapenka, V., Husband, R. J., et al. (2023). A hydrogen-enriched layer in the topmost outer core sourced from deeply subducted water. Nature Geoscience. https://doi.org/10.1038/s41561-023-01324-x

Kim, T., Wei, X., Chariton, S., Prakapenka, V. B., Ryu, Y.-J., Yang, S., and Shim, S.-H. (2023). Stability of hydrides in sub-Neptune exoplanets with thick hydrogen-rich atmospheres. Proceedings of the National Academy of Sciences, 120(52), e2309786120. https://doi.org/10.1073/pnas.2309786120

Ko, B., Greenberg, E., Prakapenka, V., Alp, E. E., Bi, W., Meng, Y., et al. (2022). Calcium dissolution in bridgmanite in the Earth's deep mantle. Nature. https://doi.org/10.1038/s41586-022-05237-4

Horn, H. W., Prakapenka, V., Chariton, S., Speziale, S., and Shim, S.-H. (2023). Reaction between hydrogen and ferrous/ferric oxides at high pressures and high temperatures - implications for sub-Neptunes and super-Earths. The Planetary Science Journal, 4(2), 30. https://doi.org/10.3847/PSJ/acab03

Wei, X., Chariton, S., Prakapenka, V. B., Fei, Y. and Shim S. (2026). Effects of hydrogen on Fe-S alloys and their implications for the martian core. JGR Planets 131, e2025JE009217.

These studies were conducted at 13IDD. 
In addition to these papers, we have two submissions from experiments performed at GSECARS, APS.

%%%%%%%%%%%%%%%%%%%%%%%%%%%%%%%%%%%%%%%%%%%%%%%%
\clearpage

\begin{figure}
\centering\includegraphics[width=0.7\textwidth]{figs/fig.jpg}
\caption{
Preliminary results for the N-H-olivine reaction. 
(a) An XRD pattern collected after heating to 3300 K at 10 GPa. 
The colored vertical bars indicate the diffraction positions of $\beta$-Si$_3$N$_4$, an Fe-N alloy phase, MgO, and residual olivine. 
The sample was laser-heated at ASU, and synchrotron XRD data were collected at Beamline 12.2.2 of the ALS. 
Diffraction peaks from olivine were detected because the X-ray beam size was too large at the ALS beamline, highlighting the importance of the high-resolution setup at 13IDD for this investigation. 
(b) Raman spectrum measured at a lower temperature region after laser heating. 
The observed peaks are consistent with N-H stretching vibrations reported in previous study \cite{chidester_Ammonia_2011} (c). 
(d) A cross-sectional SEM image \replaced{from} {corresponding to} the XRD data area in panel (a). 
The elemental maps (e, f) were collected from the boxed area revealing Fe alloy phase and silicon nitride from the reaction between N-H fluid and olivine melt. \deleted{EDS elemental distribution maps of Fe (red), Mg (green), Si (cyan), O (yellow), and N (dark blue). The spatial correlation between Si and N also supports the formation of Si$_3$N$_4$ during the reaction between olivine and H$_2$-N$_2$ fluids.
}
}\label{Olivine-N2-H2}
\end{figure}

%\begin{cbox}
%{I had to reduce the figure file size by 100 times to prevent latex crashing.  Make sure your figure file size is always less than 5 MB.}
%\end{cbox}

\clearpage

%%%%%%%%%%%%%%%%%%%%%%%%%%%%%%%%%%%%%%%%%%%%%%%%%%
%\inlinecomment{Literature references (DOIs or citations) - 2000 character limit}

\bibliographystyle{proposal2-doi-only}
\bibliography{./bibs/KimLibrary}